\documentclass[11pt]{article}
\usepackage[utf8]{inputenc}
\usepackage{amsmath,amssymb}
\usepackage{hyperref}
\title{Robust Groundwater Contamination Modeling under Distribution Shift}
\author{Anonymous}
\date{}
\begin{document}
\maketitle

\begin{abstract}
This paper studies Groundwater Contamination Modeling. Grounded in a real, citable literature \cite{ref1,ref2,ref3,ref4,ref5,ref6,ref7,ref8},
we propose a focused method and report \textbf{illustrative} experiments.
All references below are real and verifiable (OpenAlex/arXiv); the reported
numbers are simulated to illustrate intended behaviour.
\end{abstract}

\section{Introduction}
Groundwater Contamination Modeling is an active research area. This work targets a concrete gap identified from
the recent literature and contributes a method together with an evaluation protocol.

\section{Related Work (real literature)}
The following works, retrieved from public bibliographic APIs, frame our study:
\begin{itemize}
\item Meharg et al. (2002) — "Arsenic Contamination of Bangladesh Paddy Field Soils:  Implications for Rice Contribution to Arsenic Consumption", Environmental Science \& Technology (cited 1013).
\item Xue et al. (2009) — "Present limitations and future prospects of stable isotope methods for nitrate source identification in surface- and groundwater", Water Research (cited 951).
\item Hounslow (1995) — "Water Quality Data: Analysis and Interpretation", Medical Entomology and Zoology (cited 910).
\item Rodríguez‐Lado et al. (2013) — "Groundwater Arsenic Contamination Throughout China", Science (cited 905).
\item Ramakrishnaiah et al. (2008) — "Assessment of Water Quality Index for the Groundwater in Tumkur Taluk, Karnataka State, India", Journal of Chemistry (cited 898).
\item Letterman (1999) — "Water quality and treatment : a handbook of community water supplies", McGraw-Hill eBooks (cited 861).
\end{itemize}

\section{Method}
We model distribution shift explicitly and optimise a worst-case objective over an uncertainty set of plausible test distributions. We instantiate this idea for Groundwater Contamination Modeling and analyse its cost/quality trade-off.

\section{Experiments (Illustrative)}
\begin{center}
\begin{tabular}{lcc}
System & Primary Metric & Efficiency \\ \hline
Strong baseline & 0.812 & 1.00$\times$ \\
Ablation & 0.828 & 0.92$\times$ \\
\textbf{Ours} & \textbf{0.851} & \textbf{0.71$\times$}
\end{tabular}
\end{center}
\emph{Metrics simulated to illustrate the intended effect; not measured.}

\section{Conclusion}
We connected a real literature to a concrete method for Groundwater Contamination Modeling. Future work will
validate the illustrative results with real experiments.

\begin{thebibliography}{9}
\bibitem{ref1} Andrew A. Meharg, Md. Mazibur Rahman. Arsenic Contamination of Bangladesh Paddy Field Soils:  Implications for Rice Contribution to Arsenic Consumption. \emph{Environmental Science \& Technology}, 2002. DOI:10.1021/es0259842
\bibitem{ref2} Dongmei Xue, Jorin Botte, Bernard De Baets et al.. Present limitations and future prospects of stable isotope methods for nitrate source identification in surface- and groundwater. \emph{Water Research}, 2009. DOI:10.1016/j.watres.2008.12.048
\bibitem{ref3} Arthur W. Hounslow. Water Quality Data: Analysis and Interpretation. \emph{Medical Entomology and Zoology}, 1995. 
\bibitem{ref4} Luis Rodríguez‐Lado, Guifan Sun, Michael Berg et al.. Groundwater Arsenic Contamination Throughout China. \emph{Science}, 2013. DOI:10.1126/science.1237484
\bibitem{ref5} C. R. Ramakrishnaiah, C. Sadashivaiah, G. Ranganna. Assessment of Water Quality Index for the Groundwater in Tumkur Taluk, Karnataka State, India. \emph{Journal of Chemistry}, 2008. DOI:10.1155/2009/757424
\bibitem{ref6} Raymond D. Letterman. Water quality and treatment : a handbook of community water supplies. \emph{McGraw-Hill eBooks}, 1999. 
\bibitem{ref7} David I. Gustafson. Groundwater ubiquity score: A simple method for assessing pesticide leachability. \emph{Environmental Toxicology and Chemistry}, 1989. DOI:10.1002/etc.5620080411
\bibitem{ref8} Insaf S. Babiker, Mohamed A. A. Mohamed, Tetsuya Hiyama et al.. A GIS-based DRASTIC model for assessing aquifer vulnerability in Kakamigahara Heights, Gifu Prefecture, central Japan. \emph{The Science of The Total Environment}, 2005. DOI:10.1016/j.scitotenv.2004.11.005
\end{thebibliography}
\end{document}
